This thesis set out to answer two questions, one theoretical and one practical. The theoretical one has a clean answer. The practical one has a sobering answer. On a single daily price series, filtering does not buy much.

The theoretical question was how to estimate the hidden state of a noisy system step by step, and how sure we can be of that estimate. The Bayesian filtering equations of Chapter~\ref{chap:bayesian} do this for any state space model. They carry a whole distribution rather than a single number, so the uncertainty comes out along with the estimate. Whether the result is exact depends on the model. In the linear Gaussian case the Kalman filter of Chapter~\ref{chap:kalman} solves the problem exactly. Outside that case the extended Kalman filter of Chapter~\ref{chap:nonlinear} replaces the model by a straight line through the current estimate, and how good it is depends on how strong the non-linearity is.

The practical question was how well these filters predict a real price series, compared with simple baselines on data they have never seen. Letting the average return move over time changes almost nothing. The local level model, the local linear trend model and the autoregressive models all stay within about one percent of the naive forecast that predicts no change at all. The variance is a different story. The stochastic volatility model does recover the hidden volatility from the data. But recovering it does not make the next squared return easier to predict, because the volatility is buried under much more noise than signal.

These last two findings pull in different directions. What the volatility model recovers is the size of the moves, and the size of the moves is what risk means for an asset. So the filter gives a running estimate of how risky the asset is right now. That this estimate does not improve a point forecast says something about the prediction task, not about the estimate itself.

The next step forward is to model several assets together. One state can hold the hidden levels, or the hidden volatilities, of several assets at once. It then carries the correlations between them, and it can carry delayed effects too. If one asset falls and another reacts a day later, a model built on a single series cannot see it. The framework of this thesis extends to that without any change of principle, only with larger matrices.